%!TEX program = xelatex
% NOTE: For Unicode support, compile with XeLaTeX (or LuaLaTeX).
\documentclass[11pt]{article}

\usepackage{fontspec}
\usepackage{xeCJK}
%\setmainfont{CMU Serif}  % removed for portability
\IfFontExistsTF{Noto Serif CJK SC}{\setCJKmainfont{Noto Serif CJK SC}}{%
 \IfFontExistsTF{Source Han Serif SC}{\setCJKmainfont{Source Han Serif SC}}{%
  \setCJKmainfont{FandolSong-Regular}% TeX Live fallback
 }%
}

\usepackage[margin=1in]{geometry}
\usepackage{amsmath, amssymb, amsthm}
\usepackage{graphicx}
\usepackage{tikz-cd}
\usepackage{multicol}
\usepackage{natbib}
\usepackage{hyperref}
\usepackage{booktabs}

\newtheorem{definition}{Definition}
\newtheorem{assumption}{Assumption}
\newtheorem{lemma}{Lemma}
\newtheorem{theorem}{Theorem}
\newtheorem{corollary}{Corollary}
\newtheorem{proposition}{Proposition}
\newtheorem{remark}{Remark}

\setlength{\parindent}{0pt}
\setlength{\parskip}{1\baselineskip}

\title{The Mean Strikes the Hessian: Mean-Dominant Activations Impose Curvature Barriers in Large Language Models}
\author{Theory Team}
\date{\today}

\begin{document}
\maketitle

\begin{abstract}
Training Large Language Models (LLMs) often operates near the \emph{Edge of Stability} (EoS), where the classical descent-lemma condition $\eta<2/\lambda_1(\nabla^2\!L)$ is violated and yet optimization can still make progress.
Recent theory suggests that, in certain settings, EoS dynamics can be understood via a two-phase picture: iterates first approach a manifold of minima, then oscillate around it while drifting along the manifold to reduce the top Hessian eigenvalue (sharpness) \citep{arora2022understanding}.
This manuscript provides a complementary mechanism \emph{upstream of the optimizer}: we show that the empirical second moment of forward activations can contain a \emph{rank-1 mean spike} $\mu\mu^\top$ which directly enforces a lower bound on curvature in the Gauss--Newton block of an affine projection layer.
Under a mild non-saturation condition for autoregressive softmax cross-entropy, we prove
\[
\lambda_{\max}(\mathcal{H}_{\mathrm{GN}})\ =\mathcal{O}(\|\mu\|^2),
\]
up to explicitly bounded covariance-residual terms (see Theorem~\ref{theorem:learning_rate}).
Consequently, the \emph{classical} stability threshold for gradient descent scales as $\eta_{\mathrm{stab}}=\mathcal{O}(\|\mu\|^{-2})$, implying that mean-dominant activations push training toward (or deeper into) an EoS regime.
These results motivate concrete architectural interventions (mean-centering projections) and connect representation statistics to EoS-style sharpness dynamics.
\end{abstract}

\section{Introduction}

Understanding why Transformer training exhibits persistent large curvature---and why adaptive optimizers are often essential---remains a core question in modern optimization theory.
Empirically, Transformer Hessians display strong block heterogeneity, offering one explanation for why vanilla SGD underperforms and why adaptive methods (e.g., Adam) stabilize training \citep{zhang2024why}.
A second, complementary observation is the \emph{Edge of Stability} (EoS): for many deep models, deterministic GD can enter a regime where (i) sharpness rises beyond $2/\eta$;
(ii) then hovers above $2/\eta$ while slightly decreasing; and (iii) the loss oscillates per-iteration yet decreases on longer time scales \citep{cohen2021gradient}.
A recent mathematical analysis by \citet{arora2022understanding} formalizes a robust notion of \emph{stableness} and proves that, for two mechanisms (loss non-smoothness and adaptive/normalized step sizes), GD admits a two-phase description: after approaching a manifold of local minimizers, iterates overshoot across the manifold and the projection onto it follows a Riemannian gradient flow that decreases the top Hessian eigenvalue along the manifold \citep{arora2022understanding}.
This provides a principled view of \emph{why} sharpness can decrease in EoS despite local instability.
Our goal is orthogonal: rather than starting from a given loss landscape and analyzing the optimizer, we identify a \emph{representation-level} source of sharpness that can drive the system into EoS.
Empirically, activation matrices in LLMs are often spectrally anisotropic, with the leading right singular vector strongly aligned with the empirical mean activation $\mu$.
While some work attributes outlier directions to architectural artifacts such as attention sinks \citep{xiao2023efficient}, the persistence of mean-aligned spikes in models designed to mitigate such artifacts suggests a broader statistical origin.
The practical consequence is that a rank-1 mean spike in activation second moments can induce a rank-1 spike in curvature for any downstream affine projection, thereby shrinking the classical stability threshold.
Technically, we work with an \emph{exact} Gauss--Newton block for an affine map $z=Wx$ under cross-entropy loss, retaining a covariance residual term rather than assuming perfect K-FAC-style factorization.
Under mild bounded-fluctuation assumptions and a non-saturation condition, we:
\begin{enumerate}
  \item lower bound the top eigenvalue of the expected functional Hessian $\Sigma_{H_z}$;
  \item propagate the mean spike through a Kronecker-structured Gauss--Newton block.
\end{enumerate}
Finally, we discuss how these results connect to EoS theory: mean dominance provides a mechanism that \emph{sets the scale} of sharpness (and hence stableness), while EoS manifold-flow theory suggests how training might subsequently drift to reduce that sharpness.

\section{Theoretical Modeling and Mathematical Unification}

\textbf{Procedure of the theoretical derivation.}
\begin{enumerate}
  \item \textbf{Block focus:} adopt a \emph{block-heterogeneous} view of curvature and isolate a single affine projection block $z=Wx$ (e.g., a linear projection in attention/MLP/LM head).
In practice this most naturally corresponds to the lm\_head, since it is directly followed by softmax.
  \item \textbf{Assumptions:} (i) mean+residual decomposition of activations with bounded second-moment fluctuations;
(ii) non-saturation for softmax cross-entropy, yielding a positive functional-curvature floor.
  \item \textbf{Block curvature:} obtain an exact Gauss--Newton block decomposition into a Kronecker factor plus a controlled covariance residual, so the mean spike $\mu\mu^\top$ sets the sharpness scale of this \emph{specific block}.
  \item \textbf{Implications:} Under \emph{mean dominance}, the loss curvature is primarily determined by the \emph{statistics of the learned representations}, rather than by optimizer-induced dynamics.
\end{enumerate}


\subsection{Empirical Motivation and Covariance Bounds}

\textbf{Subsection Summary:}
\textit{Assumptions:} Activations can be written as $x=\mu+\tilde x$, where $\mathbb{E}[\tilde x]=0$ and second-moment fluctuations are bounded.
\textit{Main relation:} $\Sigma_{xx}=\mathbb{E}[xx^\top]=\mu\mu^\top+\Sigma_{\tilde x\tilde x}$.
\textit{Implication:} The rank-1 mean term $\mu\mu^\top$ is separated from residual covariance, so the mean contribution to sharpness can be analyzed explicitly.
Let $(x,y)\sim\mathcal{D}$ denote a token-level sample with activation $x\in\mathbb{R}^{d_{\mathrm{in}}}$ at a fixed layer.
Define $\mu := \mathbb{E}[x]$ and the centered residual $\tilde x := x-\mu$.
\begin{assumption}[Mean+residual decomposition with bounded second-moment fluctuations]
\label{assum:decomposition}
We assume
\[
x=\mu+\tilde x,\qquad \mathbb{E}[\tilde x]=0,
\]
and the second moment concentrates in spectral norm:
\[
\mathbb{E}\bigl\|xx^\top-\Sigma_{xx}\bigr\|_2^2 \le \sigma_{xx}^2,\qquad \Sigma_{xx}:=\mathbb{E}[xx^\top].
\]
\end{assumption}

Note that $\Sigma_{xx}=\mu\mu^\top+\Sigma_{\tilde x\tilde x}$ where $\Sigma_{\tilde x\tilde x}:=\mathbb{E}[\tilde x\tilde x^\top]\succeq 0$.

\subsection{PSD property and a non-trivial curvature floor for cross-entropy}

\textbf{Subsection Summary:}
\textit{Assumptions:} Non-saturation of training tokens ($c_1\le p_y\le 1-\delta$) and bounded fluctuations of $H_z$.
\textit{Main relation:} The expected functional Hessian has a positive floor: $\lambda_{\max}(\Sigma_{H_z})\ge\gamma>0$.
\textit{Implication:} In the active training regime, loss curvature remains non-degenerate, enabling a quantitative link between block sharpness and activation statistics via $\gamma$.

Consider logits $z\in\mathbb{R}^{d_{\mathrm{out}}}$, $p=\mathrm{softmax}(z)$, and cross-entropy $\ell(z,y)=-\log p_y$ with one-hot $y$.
The functional Hessian is
\[
H_z := \nabla_z^2 \ell(z,y)=\mathrm{diag}(p)-pp^\top\succeq 0.
\]
Let $\Sigma_{H_z}:=\mathbb{E}[H_z]$.

\begin{assumption}[Non-saturation and bounded functional-Hessian fluctuations]
\label{assum:nonsaturation}
There exist $\delta\in(0,1/2]$ and $c_1>0$ such that
\[
c_1 \le p_y \le 1-\delta
\quad\text{along the relevant optimization trajectory.}
\]
Moreover,
\[
\mathbb{E}\bigl\|H_z-\Sigma_{H_z}\bigr\|_2^2\le \sigma_H^2.
\]
\end{assumption}

\begin{lemma}[Curvature floor for expected functional Hessian]
\label{lemma:ce_curvature}
Under Assumption~\ref{assum:nonsaturation},
\[
\lambda_{\max}(\Sigma_{H_z}) \ge \gamma,
\qquad
\gamma:=\frac{1}{d_{\mathrm{out}}}\min\bigl(c_1(1-c_1),\ \delta(1-\delta)\bigr).
\]
\end{lemma}
\begin{proof}
For each sample, $\mathrm{tr}(H_z)=\sum_k p_k(1-p_k)=1-\|p\|_2^2\ge p_y(1-p_y)$. On $[c_1,1-\delta]$, the minimum of $x(1-x)$ occurs at endpoints, yielding $p_y(1-p_y)\ge \min(c_1(1-c_1),\delta(1-\delta))$. Taking expectation gives
\[
\mathrm{tr}(\Sigma_{H_z})\ge \min(c_1(1-c_1),\delta(1-\delta)).
\]
Since $\Sigma_{H_z}\succeq 0$, $\lambda_{\max}(\Sigma_{H_z})\ge \mathrm{tr}(\Sigma_{H_z})/d_{\mathrm{out}}$.
\end{proof}

\begin{remark}[Relaxation of the curvature barrier in the saturation regime]
\label{remark:saturation}
Assumption~\ref{assum:nonsaturation} imposes $p_y \le 1-\delta$, which yields a strictly positive curvature floor $\gamma>0$ for the expected functional Hessian during the active phase.
In the saturation regime, many tokens satisfy $p_y \to 1$, so $\delta \to 0$ and hence $\gamma \to 0$.
This only weakens the \emph{lower bound} on functional curvature; it does not by itself imply that the mean spike $\mu$ (or the rank-1 component of $\Sigma_{xx}$) disappears.
In other words, the guaranteed curvature contribution $\mathcal{O}(\gamma\|\mu\|^2)$ can vanish even if the representations remain mean-dominant.
This separation is consistent with the observation that optimization can later escape the EoS regime or tolerate larger effective step sizes once functional curvature relaxes.
\end{remark}

\subsection{Exact Gauss--Newton block with a controlled covariance residual}

\textbf{Subsection Summary:}
\textit{Assumptions:} Bounded second-moment fluctuations (Assumptions~\ref{assum:decomposition} and \ref{assum:nonsaturation}).
\textit{Main relation:} 
\[
\mathcal{H}_{\mathrm{GN}}=\Sigma_{xx}\otimes\Sigma_{H_z}+\mathcal{E}_{\mathrm{cov}},
\qquad
\|\mathcal{E}_{\mathrm{cov}}\|_2\le\epsilon:=\sigma_{xx}\sigma_H.
\]
\textit{Implication:} This gives an exact decomposition with an explicit error bound $\epsilon$, so subsequent sharpness bounds are controlled rather than heuristic.

Because $z=Wx$ is linear in $w=\mathrm{vec}(W)$, the Hessian \emph{block} w.r.t.\ $w$ equals the Gauss--Newton block:
\[
\mathcal{H}_{\mathrm{GN}} \;=\;
\mathbb{E}\bigl[J^\top H_z J\bigr],
\qquad
J = x^\top\otimes I_{d_{\mathrm{out}}}.
\]
(We emphasize: this identity holds for the $(W,W)$ block; the full-network Hessian generally differs from Gauss--Newton due to second derivatives of upstream layers.)

\begin{lemma}[Exact Kronecker decomposition with bounded residual]
\label{lemma:exact_kfac}
Define $\Sigma_{xx}=\mathbb{E}[xx^\top]$ and $\Sigma_{H_z}=\mathbb{E}[H_z]$.
Then
\begin{equation}
  \mathcal{H}_{\mathrm{GN}}
  \;=\;
  \mathbb{E}[xx^\top\otimes H_z]
  \;=\;
  \Sigma_{xx}\otimes \Sigma_{H_z} + \mathcal{E}_{\mathrm{cov}},
\end{equation}
where the covariance-residual term satisfies
\[
\|\mathcal{E}_{\mathrm{cov}}\|_2 \le \sigma_{xx}\sigma_H \;=: \epsilon.
\]
\end{lemma}

\begin{proof}
Let $A:=xx^\top-\Sigma_{xx}$ and $B:=H_z-\Sigma_{H_z}$. Then
\[
\mathcal{E}_{\mathrm{cov}}
=
\mathbb{E}[xx^\top\otimes H_z]-\Sigma_{xx}\otimes\Sigma_{H_z}
=
\mathbb{E}[A\otimes B].
\]
Using $\|A\otimes B\|_2=\|A\|_2\|B\|_2$ and Jensen,
\[
\|\mathcal{E}_{\mathrm{cov}}\|_2
=
\bigl\|\mathbb{E}[A\otimes B]\bigr\|_2
\le
\mathbb{E}\,\|A\otimes B\|_2
=
\mathbb{E}\bigl[\|A\|_2\|B\|_2\bigr]
\le
\sqrt{\mathbb{E}\|A\|_2^2\,\mathbb{E}\|B\|_2^2}.
\]
Apply Assumptions~\ref{assum:decomposition} and~\ref{assum:nonsaturation}.
\end{proof}

\subsection{Mean-induced curvature lower bound and implications for stability/EoS}

\textbf{Subsection Summary:}
\textit{Assumptions:} Mean-dominant activations and the curvature floor $\gamma>0$.
\textit{Main relation:} 
\[
\lambda_{\max}(\mathcal{H}_{\mathrm{GN}})\ \text{scales with}\ \|\mu\|^2
\quad (\text{up to }\epsilon\text{-level residual terms}).
\]
\textit{Implication:} Larger activation mean implies larger sharpness and thus a smaller classical stable step size, $\eta_{\mathrm{stab}}\sim\mathcal{O}(\|\mu\|^{-2})$.

\begin{theorem}[Mean-induced lower bound on Gauss--Newton sharpness]
\label{theorem:learning_rate}
Under Assumptions~\ref{assum:decomposition} and~\ref{assum:nonsaturation},
\begin{equation}
  \lambda_{\max}(\mathcal{H}_{\mathrm{GN}})
  \;\ge\;
\gamma\Bigl(\|\mu\|^2+\lambda_{\min}(\Sigma_{\tilde x\tilde x})\Bigr)-\epsilon,
\end{equation}
where $\Sigma_{\tilde x\tilde x}:=\mathbb{E}[\tilde x\tilde x^\top]$, $\gamma$ is from Lemma~\ref{lemma:ce_curvature}, and $\epsilon$ is from Lemma~\ref{lemma:exact_kfac}.
In particular, if $\gamma\|\mu\|^2\gg \epsilon$, then $\lambda_{\max}(\mathcal{H}_{\mathrm{GN}})=\mathcal{O}(\|\mu\|^2)$.
\end{theorem}

\begin{proof}
By Weyl's inequality,
\[
\lambda_{\max}(\mathcal{H}_{\mathrm{GN}})
\ge
\lambda_{\max}(\Sigma_{xx}\otimes\Sigma_{H_z})-\|\mathcal{E}_{\mathrm{cov}}\|_2.
\]
For PSD matrices, $\lambda_{\max}(\Sigma_{xx}\otimes\Sigma_{H_z})=\lambda_{\max}(\Sigma_{xx})\lambda_{\max}(\Sigma_{H_z})$.
Moreover, $\Sigma_{xx}=\mathbb{E}[xx^\top]=\mu\mu^\top+\Sigma_{\tilde x\tilde x}$, hence Weyl implies
\[
\lambda_{\max}(\Sigma_{xx})
\ge
\lambda_{\max}(\mu\mu^\top)+\lambda_{\min}(\Sigma_{\tilde x\tilde x})
=
\|\mu\|^2+\lambda_{\min}(\Sigma_{\tilde x\tilde x}).
\]
Combine Lemma~\ref{lemma:ce_curvature} ($\lambda_{\max}(\Sigma_{H_z})\ge\gamma$) and Lemma~\ref{lemma:exact_kfac} ($\|\mathcal{E}_{\mathrm{cov}}\|_2\le\epsilon$).
\end{proof}

\begin{remark}[On the tightness of the covariance bound and mean dominance]\label{remark:epsilon_tight}
The covariance term $\epsilon=\sigma_{xx}\sigma_H$ is exceptionally tight in sequence modeling.
Since the functional Hessian $H_z=\mathrm{diag}(p)-pp^\top$ is derived from a softmax distribution, its spectral norm is deterministically bounded by $1$ globally (i.e., $\|H_z\|_2\le 1$), which implies $\sigma_H\le 1$ as an unconditional architectural constant.
Consequently, $\epsilon\le\sigma_{xx}$ is scaled solely by residual activation fluctuations, which are empirically dwarfed by the rank-1 mean penalty $\|\mu\|^2$.
This supports the practical reasonableness of the mean-dominance condition used in Theorem~\ref{theorem:learning_rate}.
\end{remark}

\begin{remark}[From curvature to ``stableness'' and EoS onset]
\label{remark:stableness}
\citet{arora2022understanding} define the \emph{stableness} of $L$ at $(x,\eta)$ as
\[
S_L(x,\eta)
:=\eta\cdot\sup_{0\le s\le \eta}\lambda_1\!\left(\nabla^2L\bigl(x-s\nabla L(x)\bigr)\right),
\]
and call $(x,\eta)$ stable iff $S_L(x,\eta)\le 2$.
Since $\sup_{s}(\cdot)\ge \lambda_1(\nabla^2L(x))$, any lower bound on pointwise sharpness yields a lower bound on stableness.
Therefore, Theorem~\ref{theorem:learning_rate} supplies a \emph{mechanism-level} explanation for why a fixed learning rate can become (locally) unstable when $\|\mu\|$ is large: it enforces a large $\lambda_1$ and hence a large $S_L$.
In EoS theory, instability does not necessarily stop optimization; rather, the dynamics may enter a phase where sharpness is regulated and can decrease along a minimizer manifold \citep{arora2022understanding}.
\end{remark}

\begin{proposition}[Mean dominance + EoS alignment $\Rightarrow$ rank-1 updates]
\label{prop:rank1update}
Assume the Gauss--Newton block is well-approximated by $\Sigma_{xx}\otimes\Sigma_{H_z}$ with a dominant mean spike so that the top eigenvector of $\Sigma_{xx}$ is $u_\mu\approx \mu/\|\mu\|$, and let $u_H$ be the top eigenvector of $\Sigma_{H_z}$.
Then the top eigenvector of $\Sigma_{xx}\otimes\Sigma_{H_z}$ is approximately $u_\mu\otimes u_H$.
If, as predicted in EoS phase II, the gradient direction aligns with the top-eigenspace of the Hessian/GN block \citep{arora2022understanding}, then the induced update on $W$ is approximately rank-1:
\[
\Delta W \;\propto\;
-\,u_H\,u_\mu^\top \;\approx\; -\,u_H\,\frac{\mu^\top}{\|\mu\|}.
\]
Hence, mean dominance predicts that EoS-style alignment concentrates updates along the mean direction, potentially starving directions orthogonal to $\mu$.
\end{proposition}

\subsection{EoS phase II as implicit mean regularization}

\textbf{Subsection Summary:}
\textit{Assumptions:} Along the minimizer manifold, sharpness in this block is still mean-dominant (large $\rho_W$).
\textit{Main relation:} Phase-II EoS flow that decreases block sharpness also decreases (approximately) $\|\mu_W(\theta)\|^2$.
\textit{Implication:} If mean drives sharpness, then reducing sharpness and reducing the mean spike are aligned objectives, which motivates explicit mean-centering.

We briefly recall the phase-II limiting picture in \citet{arora2022understanding}. Under regularity assumptions ensuring a manifold of (local) minimizers $\Gamma$ and a smooth projection map onto it, certain finite-step mechanisms (e.g., normalized GD) admit a two-phase description: iterates first approach $\Gamma$ (phase I), then oscillate across it while the projection onto $\Gamma$ drifts to reduce sharpness (phase II).
In differential form, the projected dynamics can be written as a Riemannian gradient flow that decreases $\log\lambda_1$ along $\Gamma$:
\begin{equation}\label{eq:arora_flow}
\dot\theta(\tau) 
=
-\frac14\,\Pi_{T_{\theta(\tau)}\Gamma}\,\nabla \log \lambda_1(\theta(\tau)),
\end{equation}
where $\Pi_{T_\theta\Gamma}$ denotes orthogonal projection onto the tangent space of $\Gamma$ at $\theta$.
Consequently,
\[
\frac{d}{d\tau}\log\lambda_1(\theta(\tau))
=
\left\langle \nabla\log\lambda_1(\theta(\tau)),\,\dot\theta(\tau)\right\rangle
=
-\frac14\left\|\Pi_{T_{\theta(\tau)}\Gamma}\nabla\log\lambda_1(\theta(\tau))\right\|_2^2
\le 0,
\]
so phase II performs deterministic sharpness reduction along $\Gamma$ (until a stationary point of $\lambda_1$ on $\Gamma$ is reached).
To couple this to representation statistics, fix a layer/block $W$ and define the activation mean map
\[
\mu_W(\theta) := \mathbb{E}_{(x,y)\sim\mathcal{D}}\bigl[x_W(\theta)\bigr],
\]
and denote by $s_W(\theta):=\lambda_{\max}(\mathcal{H}_{\mathrm{GN}}^{(W)}(\theta))$ the corresponding Gauss--Newton block sharpness.
Theorem~\ref{theorem:learning_rate} gives (up to controlled residual terms) a lower bound of the form
\[
s_W(\theta)\ \gtrsim\ \gamma\|\mu_W(\theta)\|^2-\epsilon.
\]
This suggests introducing a \emph{mean-dominance ratio}
\begin{equation}\label{eq:mean_ratio}
\rho_W(\theta)
:=
\frac{\gamma\|\mu_W(\theta)\|^2}{s_W(\theta)+\epsilon}
\in(0,1],
\end{equation}
which is observable during training. Values $\rho_W(\theta)\approx 1$ mean that the mean spike explains most of the sharpness in this block.

\begin{assumption}[Mean-dominant sharpness along $\Gamma$]\label{assum:mean_dominant_flow}
Along the relevant segment of the minimizer manifold $\Gamma$, the map $\theta\mapsto\mu_W(\theta)$ is continuously differentiable and the mean-dominance ratio satisfies $\rho_W(\theta)\ge 1-\delta_0$ for some $\delta_0\in(0,1)$.
\end{assumption}

\begin{proposition}[Phase-II sharpness flow induces mean decay]\label{prop:mean_decay}
Assume that the phase-II manifold-flow mechanism applies to the sharpness proxy $s_W(\theta)$ (i.e., the relevant projected dynamics decreases $\log(s_W+\epsilon)$ along $\Gamma$), and that Assumption~\ref{assum:mean_dominant_flow} holds.
Then the projected gradient driving phase II decomposes as
\begin{equation}\label{eq:grad_relation}
\Pi_{T_\theta\Gamma}\nabla\log\bigl(s_W(\theta)+\epsilon\bigr)
=
\Pi_{T_\theta\Gamma}\nabla\log\|\mu_W(\theta)\|^2
-
\Pi_{T_\theta\Gamma}\nabla\log\rho_W(\theta).
\end{equation}
In particular, if there exists $\alpha\in[0,1)$ such that
\[
\left\|\Pi_{T_\theta\Gamma}\nabla\log\rho_W(\theta)\right\|_2
\le
\alpha\,\left\|\Pi_{T_\theta\Gamma}\nabla\log\|\mu_W(\theta)\|^2\right\|_2
\quad\text{along }\Gamma,
\]
then along the phase-II flow minimizing $\log(s_W+\epsilon)$ we have the differential inequality
\[
\frac{d}{d\tau}\log\|\mu_W(\theta(\tau))\|^2
\le
-\frac{1-\alpha}{4}\left\|\Pi_{T_{\theta(\tau)}\Gamma}\nabla\log\|\mu_W(\theta(\tau))\|^2\right\|_2^2
\le 0.
\]
\end{proposition}

\begin{proof}
By rearranging \eqref{eq:mean_ratio}, we have $s_W(\theta)+\epsilon=\gamma\|\mu_W(\theta)\|^2/\rho_W(\theta)$.
Taking logarithms and differentiating gives \eqref{eq:grad_relation}. Now consider the flow
$\dot\theta=-(1/4)\Pi_{T_\theta\Gamma}\nabla\log(s_W+\epsilon)$. Then
\[
\frac{d}{d\tau}\log\|\mu_W(\theta(\tau))\|^2
=
\left\langle \nabla\log\|\mu_W(\theta(\tau))\|^2,\,\dot\theta(\tau)\right\rangle.
\]
Substitute the flow and use \eqref{eq:grad_relation} to expand the inner product into a negative squared-norm term plus a cross term.
Control the cross term by Cauchy--Schwarz and the assumed bound on $\|\Pi_{T_\theta\Gamma}\nabla\log\rho_W(\theta)\|_2$.
\end{proof}

Proposition~\ref{prop:mean_decay} makes the following mechanism explicit: if a mean spike dominates the sharpness that phase II seeks to reduce, then phase II implicitly regularizes the activation mean (or, more generally, whatever statistic explains the dominant curvature).
This provides a theoretical justification for explicit mean-centering interventions: they align with the phase-II objective by directly setting $\mu_W\approx 0$, thereby relaxing the curvature barrier and freeing capacity for learning residual semantics.

\section{Experimental Design and Architecture Ablations}

\subsection{Empirical Validation Strategies}

To empirically test the bounds and mechanisms above, we propose three complementary strategies:

\begin{itemize}
  \item \textbf{Spectral tracking during training:} Estimate $\lambda_{\max}(\mathcal{H}_{\mathrm{GN}})$ or dominant Hessian/GN sharpness via Lanczos-style methods, and track $\|\mu\|^2$ and the ratio $\rho_W$ in \eqref{eq:mean_ratio}.
Theorem~\ref{theorem:learning_rate} predicts that when $\rho_W$ is large, stable learning rates correlate inversely with $\|\mu\|^2$.
  \item \textbf{Null-space mean-centering:} Compare a baseline Transformer to one with explicit mean-removal (e.g., $x\mapsto x-\mu$ or a learned projection that removes the top mean-aligned direction).
The theory predicts reduced sharpness and improved learning of residual semantics.
\end{itemize}

\section{Discussion and Conclusion}

Our analysis starts from an empirical observation---mean-aligned spectral spikes in activation matrices---and converts it into a curvature statement that directly constrains optimization.
The logical chain is:

\begin{itemize}
  \item \textbf{Observation $\Rightarrow$ structural assumption:} We model activations as $x=\mu+\tilde x$ with bounded second-moment fluctuations (Assumption~\ref{assum:decomposition}), isolating the mean spike $\mu\mu^\top$ as a rank-1 component in $\Sigma_{xx}$.
  \item \textbf{Loss geometry $\Rightarrow$ non-degenerate functional curvature:} A non-saturation condition for softmax cross-entropy ensures that the expected functional Hessian $\Sigma_{H_z}$ has a strictly positive spectral floor (Lemma~\ref{lemma:ce_curvature}), preventing curvature collapse during the ``active'' optimization regime.
  \item \textbf{Exact block curvature $\Rightarrow$ sharpness lower bound:} An exact Gauss--Newton decomposition with a controlled covariance residual (Lemma~\ref{lemma:exact_kfac}) lets us propagate the mean spike through a Kronecker-structured block, yielding a sharpness bound scaling as $\lambda_{\max}(\mathcal{H}_{\mathrm{GN}})=\mathcal{O}(\|\mu\|^2)$ under mean dominance (Theorem~\ref{theorem:learning_rate}).
  \item \textbf{Optimizer dynamics $\Rightarrow$ implicit mean regularization:} EoS phase-II theory predicts drift along minimizer manifolds to reduce sharpness \citep{arora2022understanding}.
In a mean-dominant regime, this drift aligns with reducing $\|\mu(\theta)\|^2$ (Proposition~\ref{prop:mean_decay}), making explicit mean-centering an intervention consistent with the phase-II objective.
\end{itemize}

\textbf{Practical significance.} The mean-dominance ratio $\rho_W$ in \eqref{eq:mean_ratio} and the bound in Theorem~\ref{theorem:learning_rate} yield testable predictions: (i) stable learning rates should correlate inversely with $\|\mu\|^2$ when $\rho_W$ is large;
(ii) explicit mean-centering should relax curvature barriers and accelerate learning of residual semantics;
and (iii) if training enters EoS phase II, $\|\mu\|^2$ should decrease (or its coupling to curvature should weaken) as sharpness is regulated.

\bibliographystyle{plainnat}
\bibliography{refs}

\end{document}